\section{General Topological Abelian Groups}

\begin{definition} A topological abelian group (TA) is a group equipped with a topology in which the group operations are continuous.
\end{definition}

\begin{remark} Sometimes we shall write $1$ for the neutral element with multiplicative notation, and sometimes we shall write $0$ for the neutral element with additive notation, depending on which one fits better the current situation. We hope this doesn't become a source for confusion.
\end{remark}

\begin{remark}
Notice that, if $G$ is an abelian group and $g \in G$, the map $x \mapsto gx$ must be a homeomorphism.
\end{remark}

\noindent We denote by $\textbf{TopAb}$ the category where objects are TA's and morphisms are continuous homomorphisms.

\begin{remark} Notice that if $G$ is a TA and $H \leq G$ a subgroup, then $H$, equipped with the induced topology, becomes a TA. Subgroups will always be assumed to be equipped with the induced topology. This is the case for any algebraic object equipped with a topology (i.e. groups, rings, monoids, etc), since the restriction of a continuous function is continuous in the induced topology.
\end{remark}


\begin{definition} Let $G$ be a TA. A fundamental system of neighborhoods around $g \in G$ is a collection of neighborhoods $\mathcal{U}$ of $g$ such that any other neighborhood $V$ of $g$ must contain at least one $U \in \mathcal{U}$. A fundamental system of neighborhoods of $1$ will be called an FSN.
\end{definition}

\begin{proposition}\label{prop:fsn_iff_topology_generated_by_translates} Let $G$ be a TA and $\mathcal{U}$ a family of open nbhds of $1$. The following are equivalent.
\begin{enumerate}
\item[(i)] $\mathcal{U}$ is a FSN.
\item[(ii)] For every $g \in G$, $g\mathcal{U}$ is a fundamental system of neighborhoods of $g$.
\end{enumerate}
\end{proposition}
\begin{remark}  In particular, this means the family $G\mathcal{U} := \{g U \;:\; g \in G, \; U \in \mathcal{U}\}$ is a basis for the topology of $G$.
\end{remark}
\begin{proof}[Proof of~\cref{prop:fsn_iff_topology_generated_by_translates}] $(\Leftarrow)$ Immediate. $(\Rightarrow)$ Let $V$ be a nbhd of some $g \in G$. Then, $g^{-1}V$ is a nbhd of $1$ and hence $U \subseteq g^{-1}V$ for some $U \in \mathcal{U}$, so $gU \subseteq V$.
\end{proof}

\begin{definition} Given an abelian group $G$ and a family $\mathcal{U}$ of subsets of $G$ such that $1 \in U$ for all $U \in \mathcal{U}$, we define the \emph{topology generated by $\mathcal{U}$} to be the topology generated by $G \mathcal{U}$.
\end{definition}

\begin{remark}By~\cref{prop:fsn_iff_topology_generated_by_translates}, for any FSN $\mathcal{U}$ of $G$, the topology generated by $\mathcal{U}$ coincides with the topology of $G$.
\end{remark}

\begin{proposition} Let $G, H$ be TA's, $\mathcal{V}$ an FSN of $H$, and $f : G \rightarrow H$ a morphism of abelian groups (not necessarily continuous). Then, the following are equivalent.
\begin{enumerate}
\item[(i)] $f$ is continuous.
\item[(ii)] For all $V \in \mathcal{V}$, $f^{-1}(V)$ contains a nbhd of $1$.
\end{enumerate}
\end{proposition}
\begin{proof} $(\Rightarrow)$ Immediate. $(\Leftarrow)$ Let $g \in G$ and $f(g) \in f(g)V$ a fundamental nbhd of $f(g)$ (by~\cref{prop:fsn_iff_topology_generated_by_translates}). The pre-image $f^{-1}(f(g) V)$ is equal to $gf^{-1}(V)$ and hence it contains $gW$ for some nbhd $W$ of $1$.
\end{proof}

\begin{proposition}\label{prop:basic_properties_about_open_sets} Let $G$ be a TA and $H \subseteq G$ a subgroup. We have the following.
\begin{enumerate}
    \item[(i)] $H$ is open $\iff G/H$ is discrete.
    \item[(ii)] $H$ is open $\implies H$ is closed.
    \item[(iii)] $H$ is closed and of finite index $\implies H$ is open.
\end{enumerate}
\end{proposition}
\begin{proof} \TODO
\end{proof}

\begin{proposition}\label{prop:topab_has_products} The category $\textbf{TopAb}$ has all (possibly infinite products). Given a family $\{G_i\}$ of TA's, their product can be realized as the product of groups equipped with the product topology, i.e. the topology generated by $\prod U_i$, where each $U_i$ is open in $G_i$ and all but finitely many $U_i$'s satisfy $U_i = G_i$.
\end{proposition}
\begin{proof} \TODO
\end{proof}


\begin{proposition}\label{prop:topab_has_coproducts} The category $\textbf{TopAb}$ has all (possibly infinite coproducts). Given a family $\{G_i\}$ of TA's, their coproduct can be realized as the coproduct of groups (i.e. the direct sum) equipped with the topology generated by $\bigoplus U_i$, where each $U_i$ is open in $G_i$.
\end{proposition}
\begin{proof} \TODO
\end{proof}

\begin{corollary} The category $\textbf{TopAb}$ has as finite biproducts.
\end{corollary}
\begin{proof} For finitely many terms $G_1, \ldots, G_n$, the product and coproduct of groups are isomorphic. Moreover, by~\cref{prop:topab_has_products} and~\cref{prop:topab_has_coproducts}, this isomorphism of groups is upgraded to an isomorphism of topological groups, since the two topologies coincide.
\end{proof}

\begin{proposition}\label{prop:injection_from_the_infinite_coproduct_to_the_product} Let $\{G_i\}$ be a family of TA's. The inclusion of groups induces an injection of TA's \[\bigoplus G_i \xhookrightarrow{} \prod G_i.\]
\end{proposition}
\begin{remark} This is \emph{not} an embedding in general: the topology on the infinite coproduct is finer than the induced topology as a subset of the product.
\end{remark}
\begin{proof}[Proof of~\cref{prop:injection_from_the_infinite_coproduct_to_the_product}] \TODO
\end{proof}

\begin{proposition}\label{prop:topab_is_almost_abelian} The category $\textbf{TopAb}$ satisfies the following.
\begin{enumerate}
\item[(i)] A morphism is a monomorphism if and only if it is injective.
\item[(ii)] A morphism is an epimorphism if and only if it is surjective.
\item[(iii)] It is pre-additive, i.e. the set of morphisms between two TA's has a group structure (the obvious one) in which composition is bilinear.
\item[(iv)] It has a zero object, namely the zero group $\{0\}$ together with the unique topology on the singleton set.
\item[(v)] It has all kernels; they are given by taking the kernel in $\textbf{Ab}$ and equipping it with the induced topology.
\item[(vi)] It has all cokernels; they are given by taking the cokernel in $\textbf{Ab}$ and equipping it with the quotient topology.
\end{enumerate}
\end{proposition}

\begin{remark} As a sidenote, if one considers the category of Hausdorff abelian groups, the cokernel changes: the cokernel of a morphism $f : G \rightarrow H$ is now the quotient by the \emph{closure of the image}, i.e. $H/\overline{f(G)}$. This goes to show that sometimes the answer is not the obvious one, and it never hurts to write these things carefully.
\end{remark}

\begin{proof}[Proof of~\cref{prop:topab_is_almost_abelian}] \TODO
\end{proof}

\begin{remark} The only remaining property for $\textbf{TopAb}$ to be upgraded to the status of abelian category is that every mono is a kernel and every epi is a cokernel. This turns out to be false. To see this, note that for a monomorphism $f : G \xhookrightarrow{} H$ to be a kernel, it must be an homeomorphism onto its image (i.e. a topological embedding). There is no reason why this should be true in general.
\end{remark}

\begin{proposition} The category $\textbf{TopAb}$ has inverse limits. More concretely, given a directed system $\{G_i\}$, the inverse limit can be realized as the inverse limit $(G, \{\pi_i : G \rightarrow G_i\})$ in $\textbf{Ab}$, equipped with the inverse limit topology, which is the \emph{coarsest} topology that makes each $\pi_i$ continuous, namely the topology generated by
$$\bigcup_i \; \{\pi_i^{-1}(U) \;:\; U \subseteq G_i \text{ open}\}.$$
\end{proposition}
\begin{proof} \TODO
\end{proof}

\begin{proposition} The category $\textbf{TopAb}$ has direct limits. More concretely, given a directed system $\{G_i\}$, the direct limit can be realized as the direct limit $(G, \{\tau_i : G_i \rightarrow G\})$ in $\textbf{Ab}$, equipped with the direct limit topology, which is the \emph{finest} topology which makes each $\tau_i$ continuous, namely the topology generated by $$\bigcup_i \; \{\tau_i(U) \;:\; U \subseteq G_i \text{ open}\}.$$
\end{proposition}
\begin{proof} \TODO
\end{proof}


\subsection{Linear topologies}



\subsection{Completions}

\begin{definition} Let $G$ be an LTA and $\mathcal{U} = \{U_i\}$ an FSG. We define the $\mathcal{U}$-completion of $G$ as
\[G_{\mathcal{U}}^\land := \lim_{\leftarrow} G/U_i.\]
\end{definition}

\noindent We write $\pi_U = \pi_U^{\mathcal{U}}$ for a canonical projection $G_{\mathcal{U}}^\land \rightarrow G/U$. We have a natural map $c_\mathcal{U} : G \rightarrow G_\mathcal{U}^\land$ (given by $g \mapsto (\text{class of }g\text{ in } G/U_i)_i$). This map is continuous. Indeed, by the definition of the topology on the inverse limit, this map is continuous if and only if it is continuous after composing with the projections, which yields the trivially continuous map $G \rightarrow G/U_i$.

\begin{proposition}[All completions are the same]\label{prop:completions_are_all_the_same} Let $G$ be an LTA. For every pair of FSG's $\mathcal{U}, \mathcal{V}$, there exists a natural isomorphism $\phi_{\mathcal{U}, \mathcal{V}} : G_\mathcal{U}^\land \cong G_\mathcal{V}^\land$. Moreover, these maps satisfy the following properties.
\begin{enumerate}
\item[(i)] For all FSG's $\mathcal{U}, \mathcal{V}$ and $U \in \mathcal{U} \cap \mathcal{V}$, we have $\pi_{U}^{\mathcal{V}} \circ \phi_{\mathcal{U}, \mathcal{V}} = \pi_U^\mathcal{U}$.
\item[(ii)] For all FSG's $\mathcal{U}, \mathcal{V}, \mathcal{W}$, we have $\phi_{\mathcal{U}, \mathcal{W}} = \phi_{\mathcal{V}, \mathcal{W}} \circ \phi_{\mathcal{U}, \mathcal{V}}$.
\item[(iii)] For all FSG's $\mathcal{U}, \mathcal{V}$, we have $\phi_{\mathcal{U}, \mathcal{V}} \circ c_\mathcal{U} = c_{\mathcal{V}}$.
\end{enumerate}
\end{proposition}
\begin{proof} \TODO
\end{proof}

\noindent \cref{prop:completions_are_all_the_same} allows us to talk about \emph{the} completion $G^\land$, since all completions $G_{\mathcal{U}}^\land$ are canonically identified. Perhaps a more concrete way to view $G^\land$ is as the completion over the FSG of \emph{all} neighborhoods of $1$. In fact, we could have skipped all of these identification and just that as the definition of completion, but it is very useful to be able to "restrict" ourselves to smaller FSG's. Notice also that condition (iii) ensures that we have a canonical map $c : G \rightarrow G^\land$. Moreover, note that for any open subgroup $U \subseteq G$ we have a canonical morphism $G^\land \rightarrow G/U$.

\begin{definition}\label{def:c_map_and_cauchy_sequences} Let $G$ be a TA. We have the following.
\begin{enumerate}
    \item[(i)] $G$ is said to be Hausdorff if the map $c : G \rightarrow G^\land$ is injective.
    \item[(ii)] $G$ is said to be complete if the map $c : G \rightarrow G^\land$ is surjective.
\end{enumerate}
\end{definition}

\noindent Below we show that the above definitions agrees with the usual notion of Hausdorff for topological spaces and complete for metric spaces.

\begin{proposition} Let $G$ be an LTA. We have the following.
\begin{enumerate}
    \item[(i)] $G$ is Hausdorff if and only if points can be separated, i.e. for any two points $g, h \in G$ there are open nbhds $g \in U, h \in V$ such that $U \cap V = \varnothing$.
    \item[(ii)] If the topology on $G$ is metrizable, then $G$ is complete if and only if every cauchy sequence converges.
\end{enumerate}
\end{proposition}
\begin{proof} \TODO
\end{proof}

\begin{proposition} We have the following.
\begin{enumerate}
    \item[(i)] The completion $G^\land$ is always complete and Hausdorff.
    \item[(ii)] The morphism $c : G \rightarrow G^\land$ is universal in the following sense: for any morphism $f : G \rightarrow H$, with $H$ a complete and Hausdorff LTA, there is a unique morphism $\phi : G^\land \rightarrow H$ such that $\phi \circ c = f$.
\end{enumerate}
\end{proposition}
\begin{proof} \TODO
\end{proof}

\begin{proposition} Let $G$ be an LTA and $U \subseteq G$ an open nbhd of $1$. Then, we have a natural isomorphism $$U^\land \cong {\rm Ker}(G^\land \rightarrow G/U)$$
given by, from right to left, $(g_W)_W \mapsto (g_{W})_{W \subseteq U}$.
\end{proposition}
\begin{proof} \TODO
\end{proof}

\begin{remark} Note that, by~\cref{prop:basic_properties_about_open_sets}, since $U \subseteq G$ is open, the quotient $G/U$ is discrete, and hence ${\rm Ker}(G^\land \rightarrow G/U) = \pi_U^{-1}(1)$ is open. Therefore $U^\land$ can be canonically seen as an open subgroup of $G^\land$. Moreover, under this isomorphism the projections of $U^\land$ become clear: and element $\widehat{g} \in U^\land$ is an element $\widehat{g} \in G^\land$ whose projection to $G/U$ is zero; this means that the projection to $G/W$ for any $W \subseteq U$ must, in fact, factor through $U/W$.
\end{remark}

\begin{proposition} Let $G$ be an LTA. If $\mathcal{U} = \{U_i\}$ is an FSG of $G$, then $\mathcal{U}^\land := \{U_i^\land\}$ is an FSG for $G^\land$.
\end{proposition}
\begin{proof} \TODO
\end{proof}

\begin{proposition} Completion commutes with arbitrary products. More concretely, if $\{G_i\}$ is a family of LTA's, there is a natural isomorphism $$\prod G_i^\land \cong ( \prod G_i )^\land.$$
\end{proposition}
\begin{proof} \TODO
\end{proof}

\noindent Since $\textbf{TopAb}$ has all finite biproducts, completion must commute with finite coproducts as well. But it turns out that it does not commute with infinite coproducts. The best we can do is the following.

\begin{proposition} We have cannonical injections
\[\bigoplus G_i^\land \xhookrightarrow{} (\bigoplus G_i)^\land \xhookrightarrow{} \prod G_i^\land\]
whose composition is the injection from~\cref{prop:injection_from_the_infinite_coproduct_to_the_product}.
\end{proposition}
\begin{proof} \TODO
\end{proof}

\begin{remark} An example of why the first map is not always surjective can be given as follows. Let $p$ be a prime number. Consider $\mathbb{Z}$ with the $p$-adic topology, i.e. that generated by the FSG $\{p^k \mathbb{Z} \;:\; k \geq 0\}$. The completion $\mathbb{Z}^\land$ becomes the $p$-adic integers $\mathbb{Z}_p$. \TODO Finish this example
\end{remark}

